% Provenance (never printed): see sources/outline_sources.md (Year 2: tholeiitic, calc-alkaline
% and alkaline series; the TAS diagram; settings; isotopic signatures of mantle and crust).
% Classification from BGS RR 99-06 (after Streckeisen 1976, Le Bas et al. 1986, Le Maitre et
% al. 1989); series and data from Vermeesch and Pease (2021) and the GeoplotR package;
% curves and counts from figdata/bachelor-2/series.py; exercise and problem values in the
% Book 4 scratch calculators calc/ch09.py and calc/ch09pb.py.
\chapter{Magmatische reeksen en tektonische omgevingen}\label{ch:b2:magmatic-series}

Geef een petroloog een chemische analyse en geen kaart, en ze kan vaak zeggen of de
lava uitvloeide bij een rug, op een oceaaneiland of boven een subducerende plaat. De lava's
van IJsland en die van de Cascadevulkanen in het westen van Noord-Amerika lopen allebei
van basalt tot rhyoliet; met honderden geanalyseerd, volgen ze twee verschillende paden terwijl
ze evolueren, het ene steeds rijker aan ijzer, het andere niet. Dit hoofdstuk geeft eerst
stollingsgesteenten een naam, naar hun mineralen wanneer die zichtbaar zijn en naar hun
chemie wanneer dat niet zo is; daarna groepeert het ze in reeksen, families van gesteenten
die door differentiatie verwant zijn, en verbindt het elke reeks met een tektonische omgeving, met de
diagrammen waarmee petrologen een omgeving uit een analyse aflezen, en hun grenzen.
% ledger: vp21:series

\begin{recall}
\textbf{Uit Boek~2}: basalt, andesiet, rhyoliet, gabbro en graniet;
felsische en mafische gesteenten; vulkanische en plutonische texturen. \textbf{Uit
Boek~3}: mid-oceanische ruggen, subductiezones en vulkanische bogen, hotspots en grote
stollingsprovincies, continentale riften. \textbf{Uit \cref{ch:b2:elements-in-the-earth}}:
MORB en OIB, op de \omterm{def:b2:elements-in-the-earth:bse}{primitieve mantel} genormaliseerde (spin)diagrammen, compatibele en
\omterm{def:b2:elements-in-the-earth:partition}{incompatibele elementen}. \textbf{Uit \cref{ch:b2:isotope-geochemistry}}: de \omterm{def:b2:isotope-geochemistry:reservoirs}{verarmde
mantel} en de verrijkte isotopenreservoirs. \textbf{Uit
\cref{ch:b2:magmatic-differentiation}}: \omterm{def:b2:magmatic-differentiation:crystallisation}{fractionele kristallisatie}, de \omterm{def:b2:magmatic-differentiation:lld}{afdalingslijn van de
vloeistof}, \omterm{def:b2:magmatic-differentiation:harker}{Harkerdiagrammen} en de wet van Rayleigh.
\end{recall}

\section{Stollingsgesteenten benoemen}\label{sec:b2:magmatic-series:naming}

Een grofkorrelig gesteente toont zijn mineralen en wordt naar hen benoemd. Een lava is daarvoor vaak te
fijnkorrelig of te glasachtig, en wordt naar zijn chemische analyse benoemd. Beide
schema's zijn internationale afspraken: ze beschrijven, en zeggen niets over hoe
het gesteente is ontstaan.
% ledger: bgs:tas-use

\begin{definition}[Modale samenstelling; QAPF-diagram]\label{def:b2:magmatic-series:qapf}
De \emph{modale samenstelling}\index{modale samenstelling} van een gesteente is het
volumeaandeel van elk van zijn mineralen. Het \emph{QAPF-diagram}\index{QAPF-diagram} benoemt een
grofkorrelig stollingsgesteente met minder dan \qty{90}{\percent} mafische mineralen naar
vier modale componenten, herrekend tot een som van 100: kwarts (Q), alkaliveldspaat
(A), plagioklaas (P) en de foïden (F). Omdat kwarts en foïden niet
samen kristalliseren, valt elk gesteente in een van twee driehoeken, QAP of FAP, die
langs hun gemeenschappelijke zijde AP aan elkaar grenzen.
\end{definition}
% ledger: bgs:qapf

\begin{omfigure}
\centering
\begin{tikzpicture}[font=\scriptsize, x=1cm, y=1cm]
\def\L{6.4}
\pgfmathsetmacro{\H}{\L*0.8660254}
% upper point (q, p): q = Q fraction, p = P/(A+P)
\newcommand{\qapfup}[2]{({(#1)*\L/2 + (#2)*(1-(#1))*\L}, {(#1)*\H})}
\newcommand{\dn}[2]{({(#1)*\L/2 + (#2)*(1-(#1))*\L}, {-(#1)*\H})}
\draw[thick] (0,0) -- (\L/2,\H) -- (\L,0) -- (\L/2,-\H) -- cycle;
\draw (0,0) -- (\L,0);
\foreach \q in {0.9,0.6,0.2,0.05} \draw \qapfup{\q}{0} -- \qapfup{\q}{1};
\foreach \q in {0.1,0.6} \draw \dn{\q}{0} -- \dn{\q}{1};
\foreach \p in {0.1,0.65,0.9} \draw \qapfup{0.6}{\p} -- \qapfup{0}{\p};
\draw[dashed] \qapfup{0.6}{0.35} -- \qapfup{0.2}{0.35};
\draw \qapfup{0.2}{0.35} -- \qapfup{0}{0.35};
\foreach \p in {0.1,0.35,0.65,0.9} \draw \dn{0}{\p} -- \dn{0.1}{\p};
\foreach \p in {0.1,0.5,0.9} \draw \dn{0.1}{\p} -- \dn{0.6}{\p};
\node[anchor=south] at (\L/2,\H) {\normalsize Q};
\node[anchor=east] at (0,0) {\normalsize A};
\node[anchor=west] at (\L,0) {\normalsize P};
\node[anchor=north] at (\L/2,-\H) {\normalsize F};
\node at \qapfup{0.94}{0.5} {1a};
\node at \qapfup{0.75}{0.5} {1b};
\node[font=\tiny] at \qapfup{0.4}{0.05} {2};
\node at \qapfup{0.4}{0.225} {3a};
\node at \qapfup{0.4}{0.5} {3b};
\node at \qapfup{0.4}{0.775} {4};
\node[font=\tiny] at \qapfup{0.4}{0.95} {5};
\foreach \p/\n in {0.225/7*,0.5/8*,0.775/9*} \node at \qapfup{0.125}{\p} {\n};
\node[font=\tiny] at \qapfup{0.125}{0.05} {6*};
\node[font=\tiny] at \qapfup{0.125}{0.95} {10*};
\foreach \p/\n in {0.225/7,0.5/8,0.775/9} \node at \qapfup{0.025}{\p} {\n};
\foreach \p/\n in {0.225/7',0.5/8',0.775/9'} \node at \dn{0.05}{\p} {\n};
\node[font=\tiny] at \dn{0.3}{0.05} {11};
\node at \dn{0.35}{0.3} {12};
\node at \dn{0.35}{0.7} {13};
\node[font=\tiny] at \dn{0.3}{0.95} {14};
\node at \dn{0.75}{0.5} {15};
\node[anchor=east] at \qapfup{0.9}{0} {90};
\node[anchor=east] at \qapfup{0.6}{0} {60};
\node[anchor=east] at \qapfup{0.2}{0} {20};
\node[anchor=east] at \dn{0.1}{0} {10};
\node[anchor=east] at \dn{0.6}{0} {60};
\foreach \p/\n in {0.1/10,0.35/35,0.65/65,0.9/90} \node[anchor=south, fill=white, inner sep=0.5pt] at \qapfup{0.6}{\p} {\tiny\n};
\node[anchor=west, align=left, font=\scriptsize] at (\L+0.9,1.6) {Bovenste driehoek: \% Q\\in Q + A + P;\\
  lijnen met constante\\$100\,P/(A+P)$\\bij 10, 35, 65, 90};
\node[anchor=west, align=left, font=\scriptsize] at (\L+0.9,-1.7) {Onderste driehoek: \% F\\in F + A + P;\\
  $100\,P/(A+P)$\\bij 10, 50, 90};
\end{tikzpicture}
\par\smallskip{\small De QAPF-classificatie van grofkorrelige stollingsgesteenten. 1a
kwartsoliet, 1b kwartsrijk granitoïde, 2 alkaliveldspaatgraniet, 3a syenograniet en
3b monzograniet (samen graniet), 4 granodioriet, 5 tonaliet; 6*--10* de
kwartshoudende varianten van 6--10: 6 alkaliveldspaatsyeniet, 7 syeniet, 8 monzoniet, 9 monzodioriet
of monzogabbro, 10 dioriet, gabbro of anorthosiet; 6'--10' hun foïdhoudende
varianten; 11 foïdsyeniet, 12 foïdmonzosyeniet, 13 foïdmonzodioriet of
foïdmonzogabbro, 14 foïddioriet of foïdgabbro, 15 foïdoliet. De smalle velden 6, 6', 10 en 10'
langs de randen zijn niet genummerd. Naar het schema van de British Geological Survey (gebaseerd
op Streckeisen 1976).}
\end{omfigure}
% ledger: bgs:qapf

Veld 10 draagt drie namen. Het \omterm{def:b2:magmatic-series:qapf}{QAPF-diagram} onderscheidt dioriet, gabbro en
anorthosiet niet, noch de gabbroïde gesteenten onderling; verdere diagrammen, gebaseerd op het
plagioklaas en de mafische mineralen, doen dat wel.

\begin{definition}[Totaal-alkali--silicadiagram]\label{def:b2:magmatic-series:tas}
Het \emph{totaal-alkali--silicadiagram}\index{totaal-alkali--silicadiagram} (TAS-diagram)
benoemt een vulkanisch gesteente waarvan de modus niet kan worden gemeten, naar zijn
massapercentages \ce{Na2O + K2O} tegen \ce{SiO2}, nadat de analyse is
herrekend tot \qty{100}{\percent} zonder water en koolstofdioxide. Het is verdeeld
in vijftien velden door rechte lijnen waarvan de snijpunten vaste coördinaten hebben
(de figuur na de volgende alinea).
\end{definition}
% ledger: bgs:tas, bgs:tas-use

Het diagram heeft drie groepen velden. Langs de onderkant, bij lage alkaligehalten, liggen de
namen van de gewone reeksen, pikrobasalt en basalt tot rhyoliet, gescheiden bij 45, 52,
57 en \qty{63}{\percent} silica en, voor rhyoliet, door een hellende lijn. Daarboven bevatten de S-velden en trachiet gesteenten die bij hetzelfde silicagehalte rijker zijn aan
alkaliën; de U-velden, fonoliet en foidiet bevatten gesteenten die zo rijk zijn aan
alkaliën, en zo arm aan silica, dat ze foïden kristalliseren. Enkele velden worden
door een tweede criterium onderverdeeld: een trachybasalt heet hawaiiet als $\ce{Na2O} - 2
\ge \ce{K2O}$, anders kalitrachybasalt; een tefriet met meer dan
\qty{10}{\percent} olivijn in zijn norm (hieronder) is een basaniet; een rhyoliet of een trachiet waarvan
de moleculaire $(\ce{Na2O + K2O})/\ce{Al2O3}$ groter is dan 1, heet peralkalisch. Het schema
vraagt om analyses met minder dan \qty{2}{\percent} water en \qty{0.5}{\percent}
koolstofdioxide: verwering voegt water toe en verplaatst de alkaliën, en een \omterm{def:b2:magmatic-differentiation:cumulate}{cumulaat} is geen
vloeistof.
% ledger: bgs:tas, bgs:tas-use

\begin{omfigure}
\centering
\begin{tikzpicture}
\begin{axis}[omrecord, width=0.86\linewidth, height=8.6cm,
  xmin=37, xmax=80, ymin=0, ymax=16, xtick={37,41,...,77}, ytick={0,2,...,16},
  xlabel={\ce{SiO2} (\unit{\percent})}, ylabel={\ce{Na2O + K2O} (\unit{\percent})},
  legend style={font=\scriptsize, at={(0.98,0.98)}, anchor=north east},
  tick label style={font=\scriptsize}, label style={font=\small}]
\addplot[only marks, mark=*, mark size=0.55pt, blue!70!black, opacity=0.6] table[x=SiO2, y=alk] {figdata/out/bachelor-2-series-cath-th.dat};
\addlegendentry{IJsland}
\addplot[only marks, mark=*, mark size=0.55pt, red!75!black, opacity=0.6] table[x=SiO2, y=alk] {figdata/out/bachelor-2-series-cath-ca.dat};
\addlegendentry{Cascaden}
\foreach \f in {Pc,B,O1,O2,O3,R,S1,S2,S3,T,U1,U2,U3,Ph,F}
  \addplot[black, thin, forget plot] table[x=SiO2, y=alk] {figdata/out/bachelor-2-series-tas-\f.dat};
\node[font=\tiny] at (axis cs:43,1.6) {Pc};
\node[font=\tiny, fill=white, inner sep=0.5pt] at (axis cs:50.5,0.8) {B};
\node[font=\tiny] at (axis cs:54.5,2.2) {O1};
\node[font=\tiny] at (axis cs:60,2.2) {O2};
\node[font=\tiny] at (axis cs:66.5,2.2) {O3};
\node[font=\tiny] at (axis cs:74.5,11) {R};
\node[font=\tiny, fill=white, inner sep=0.5pt] at (axis cs:49,5.8) {S1};
\node[font=\tiny] at (axis cs:52.5,7.4) {S2};
\node[font=\tiny] at (axis cs:57.5,8.9) {S3};
\node[font=\tiny] at (axis cs:64.5,12) {T};
\node[font=\tiny] at (axis cs:44,6.6) {U1};
\node[font=\tiny] at (axis cs:48.7,9.4) {U2};
\node[font=\tiny] at (axis cs:52.5,11.8) {U3};
\node[font=\tiny] at (axis cs:57,14.6) {Ph};
\node[font=\tiny] at (axis cs:39,10) {F};
\end{axis}
\end{tikzpicture}
\par\smallskip{\small Het TAS-diagram (naar Le Bas en anderen, 1986, zoals gegeven door de
British Geological Survey) met 456 lava's en intrusiegesteenten uit IJsland (blauw) en
580 uit de Cascaden (rood), oxiden zoals geanalyseerd. Pc pikrobasalt, B basalt, O1 basaltisch
andesiet, O2 andesiet, O3 daciet, R rhyoliet; S1 trachybasalt, S2 basaltisch
trachyandesiet, S3 trachyandesiet, T trachiet of trachydaciet; U1 tefriet of
basaniet, U2 fonotefriet, U3 tefrifonoliet, Ph fonoliet, F foidiet.}
\end{omfigure}
% ledger: bgs:tas, geoplotr:tas, data:b4-cath


\begin{definition}[Silicaverzadiging]\label{def:b2:magmatic-series:saturation}
Een \omterm{def:b2:magma-emplacement:magma}{magma} of gesteente is \emph{oververzadigd aan silica}\index{silicaverzadiging} als het
meer silica heeft dan zijn andere oxiden in silicaatmineralen kunnen binden, zodat
kwarts (of een ander silicamineraal) verschijnt; \emph{onderverzadigd aan silica} als het te
weinig heeft om al zijn alkaliën in veldspaat om te zetten, zodat in plaats daarvan een foïde (nefelien,
leuciet) verschijnt, nooit samen met kwarts; en \emph{verzadigd aan silica}
tussen die twee in.
\end{definition}

Voor een lava waarvan de mineralen niet te tellen zijn, wordt de verzadiging beoordeeld op een \emph{norm}: een
standaardberekening die de oxiden van de analyse in een vaste volgorde over ideale mineralen
verdeelt, alkaliën eerst in veldspaat, en het silicaoverschot als
kwarts opgeeft of het tekort als foïde. De gebruikelijke, de CIPW-norm, scheidt de
twee soorten basalt: een basalt met nefelien in zijn norm is een alkalibasalt; een
zonder is een subalkalische basalt, en daaronder hebben de tholeïtische basalten
hypersteen (orthopyroxeen) in de norm en, vaak, calciumarm pyroxeen in de modus.
\Cref{exo:b2:magmatic-series:11} voert er een vereenvoudigde versie van uit.
% ledger: bgs:tas-use, bgs:series-terms

\section{Magmatische reeksen}\label{sec:b2:magmatic-series:series}

Lava's van één vulkaan, of één provincie, zijn zelden gelijk, maar ze zijn verwant: ze
ontstaan uit vergelijkbare oudermagma's door de differentiatie van
\cref{ch:b2:magmatic-differentiation}, en ze tekenen samen een afdalingslijn.

\begin{definition}[Magmatische reeks]\label{def:b2:magmatic-series:series}
Een \emph{magmatische reeks}\index{magmatische reeks} is een reeks stollingsgesteenten van één
magmatische cyclus of provincie die door differentiatie uit vergelijkbare oudermagma's verwant zijn
en kenmerkende chemische en mineralogische eigenschappen delen. Gewoonlijk worden er drie
onderscheiden. De \emph{tholeïtische reeks}\index{tholeïtische reeks} begint bij
basalten die verzadigd of licht oververzadigd zijn aan silica, met hypersteen in de norm,
en vertoont aan het begin van haar differentiatie een duidelijke verrijking in ijzer. De
\emph{kalkalkalische reeks}\index{kalkalkalische reeks} loopt van basalt via
andesiet naar daciet en rhyoliet, met weinig of geen verrijking in ijzer, en met
waterhoudende mineralen (hoornblende, biotiet) in haar geëvolueerde leden. De \emph{alkalische
reeks}\index{alkalische reeks} begint bij basalten met nefelien in de norm en evolueert
via de S- en U-velden van het TAS-diagram naar trachiet en fonoliet.
\end{definition}
% ledger: bgs:series-terms, bgs:tas-use

De namen beschrijven reeksen, geen afzonderlijke gesteenten: een losse basalt kan tot elk van
hen behoren, en dezelfde namen zijn met verschillende betekenissen gebruikt, dus een petroloog
moet zeggen wat ze ermee bedoelt. Het onderscheid tussen de tholeïtische en de
\omterm{def:b2:magmatic-series:series}{kalkalkalische reeks} berust op ijzer, en het diagram dat dat het best toont, zet drie
groepen oxiden uit.
% ledger: bgs:series-terms

\begin{definition}[AFM-diagram]\label{def:b2:magmatic-series:afm}
Een \emph{AFM-diagram}\index{AFM-diagram} is een driehoeksdiagram met als hoekpunten
$A = \ce{Na2O + K2O}$, $F = \ce{FeO_T}$ (al het ijzer als \ce{FeO}) en $M = \ce{MgO}$, in
massaprocent herrekend zodat $A + F + M = 100$. Een lijn erop scheidt het
tholeïtische veld, aan de kant van F, van het kalkalkalische veld.
\end{definition}
% ledger: vp21:series, geoplotr:afm

Een differentiërend basalt verliest magnesium sneller dan wat ook, omdat olivijn en
pyroxeen rijker aan magnesium zijn dan de vloeistof, en houdt zijn alkaliën, die
incompatibel zijn: zijn afdalingslijn beweegt van M weg, naar A. Wat er met
ijzer gebeurt, bepaalt de reeks. In een tholeïtisch \omterm{def:b2:magma-emplacement:magma}{magma} vertrekt ijzer alleen met de
magnesiumrijke ijzer(II)-silicaten, trager dan magnesium, zodat de verhouding
$F/M$ stijgt en de trend eerst naar F klimt voordat hij naar A afbuigt: de
\emph{Fennertrend}. In een kalkalkalisch \omterm{def:b2:magma-emplacement:magma}{magma}, sterker geoxideerd, vertrekt ijzer ook in
magnetiet, dat ijzer(III) bevat, en blijft $F/M$ bijna constant terwijl de vloeistof
rechtstreeks naar A beweegt: de \emph{Bowentrend}. De figuur hieronder toont beide met
echte gesteenten, de tholeïeten van IJsland en de kalkalkalische lava's van de Cascaden. De
scheidingslijn, in 1971 op het oog getrokken, scheidt ze goed waar ze basaltisch zijn; bij
hoge alkaligehalten komen de twee trends samen en zegt de lijn weinig.
% ledger: vp21:series, geoplotr:afm, data:b4-cath

\begin{omfigure}
\centering
\begin{tikzpicture}
\begin{axis}[name=afm, width=0.5\linewidth, height=0.5\linewidth, axis equal image,
  hide axis, xmin=-0.04, xmax=1.04, ymin=-0.06, ymax=0.92, clip=false]
\addplot[only marks, mark=*, mark size=0.5pt, blue!70!black, opacity=0.6] table[x=afm_x, y=afm_y] {figdata/out/bachelor-2-series-cath-th.dat};
\addplot[only marks, mark=*, mark size=0.5pt, red!75!black, opacity=0.6] table[x=afm_x, y=afm_y] {figdata/out/bachelor-2-series-cath-ca.dat};
\addplot[very thick, black] table[x=x, y=y] {figdata/out/bachelor-2-series-afm-line.dat};
\draw[thick] (axis cs:0,0) -- (axis cs:1,0) -- (axis cs:0.5,0.8660254) -- cycle;
\node[anchor=south] at (axis cs:0.5,0.866) {F};
\node[anchor=north east] at (axis cs:0,0) {A};
\node[anchor=north west] at (axis cs:1,0) {M};
\node[font=\scriptsize, fill=white, inner sep=1pt] at (axis cs:0.52,0.62) {tholeïtisch};
\node[font=\scriptsize, fill=white, inner sep=1pt] at (axis cs:0.33,0.12) {kalkalkalisch};
\end{axis}
\begin{axis}[omrecord, at={(afm.east)}, anchor=west, xshift=1.3cm,
  width=0.46\linewidth, height=0.46\linewidth,
  xmin=-2.5, xmax=2.2, ymin=-4.5, ymax=1.5,
  xlabel={$\ln(A/F)$}, ylabel={$\ln(M/F)$},
  legend style={font=\scriptsize, at={(0.98,0.98)}, anchor=north east},
  tick label style={font=\scriptsize}, label style={font=\small}]
\addplot[only marks, mark=*, mark size=0.5pt, blue!70!black, opacity=0.6] table[x=lnAF, y=lnMF] {figdata/out/bachelor-2-series-lr-th.dat};
\addlegendentry{IJsland}
\addplot[only marks, mark=*, mark size=0.5pt, red!75!black, opacity=0.6] table[x=lnAF, y=lnMF] {figdata/out/bachelor-2-series-lr-ca.dat};
\addlegendentry{Cascaden}
\addplot[thick, blue!50!black, forget plot] table[x=x, y=y_th] {figdata/out/bachelor-2-series-lr-fits.dat};
\addplot[thick, red!50!black, forget plot] table[x=x, y=y_ca] {figdata/out/bachelor-2-series-lr-fits.dat};
\end{axis}
\end{tikzpicture}
\par\smallskip{\small Links: het \omterm{def:b2:magmatic-series:afm}{AFM-diagram} van de IJslandse (blauw) en Cascadische (rood)
gesteenten van de TAS-figuur, met de scheidingslijn van Irvine en Baragar (1971). Rechts:
dezelfde gesteenten in log-ratiocoördinaten, waar elke reeks een bijna rechte
band wordt; de lijnen zijn orthogonale regressies, met hellingen $-1.10$ (IJsland) en $-0.72$
(Cascaden), die elkaar snijden in $(-3.84, 2.12)$, zoals Vermeesch en Pease (2021) vonden.}
\end{omfigure}
% ledger: geoplotr:afm, data:b4-cath, vp21:model

Van de IJslandse gesteenten binnen het bereik van de scheidingslijn vallen er 351 aan de
tholeïtische kant en 23 aan de andere; van de Cascadische gesteenten 417 aan de kalkalkalische
kant en 57 aan de andere. Een kromme is een onhandige grens, en een eenvoudig model verklaart
de vorm van de trends. Stel dat een \omterm{def:b2:magma-emplacement:magma}{magma} zijn alkaliën, ijzer en magnesium verliest
met snelheden evenredig met de hoeveelheden die het bevat, $\mathrm dX/\mathrm dt = -\lambda_X X$.

\begin{proposition}[Log-ratiotrends]\label{prop:b2:magmatic-series:logratio}
Als $A$, $F$ en $M$ exponentieel afnemen met constanten $\lambda_A$, $\lambda_F$ en
$\lambda_M$, beweegt het \omterm{def:b2:magma-emplacement:magma}{magma} langs een rechte lijn in het vlak van $\ln(A/F)$ en
$\ln(M/F)$, met helling $(\lambda_F - \lambda_M)/(\lambda_F - \lambda_A)$, terwijl het
langs een kromme beweegt in de AFM-driehoek.
\end{proposition}

\begin{proof}
Met $X = X_0 e^{-\lambda_X t}$ geldt $\ln(A/F) = \ln(A_0/F_0) + (\lambda_F - \lambda_A)t$
en $\ln(M/F) = \ln(M_0/F_0) + (\lambda_F - \lambda_M)t$: beide zijn lineair in $t$, dus
de een is lineair in de ander, met de genoemde helling. De verhoudingen $A/F$ en $M/F$ leggen het
punt in de driehoek vast, waarvan de coördinaten $A/(A + F + M)$, \dots\ niet lineair zijn in $t$.
\end{proof}

De steilere helling van de IJslandse gesteenten betekent een kleinere $\lambda_F$: ijzer verlaat
het tholeïtische \omterm{def:b2:magma-emplacement:magma}{magma} trager (\cref{exo:b2:magmatic-series:10}). Vermeesch en Pease toonden aan dat de twee
banden van de figuur hierboven in één punt samenkomen, dat ze opvatten als een hypothetische gemeenschappelijke
bron.
% ledger: vp21:model

De tholeïtische en de \omterm{def:b2:magmatic-series:series}{kalkalkalische reeks} verschillen ook in hun kali. Het
\ce{K2O}--\ce{SiO2}-diagram verdeelt basalten, andesieten, dacieten en rhyolieten met
meer dan \qty{48}{\percent} silica in laag-K-, middel-K- en hoog-K-typen met twee
rechte lijnen (de figuur hieronder). De meeste gesteenten van beide reeksen hier zijn middel-K: het
diagram zegt hoeveel kali een reeks meedraagt, niet welke reeks het is.
% ledger: bgs:k2o-sio2, data:b4-cath

\begin{omfigure}
\centering
\begin{tikzpicture}
\begin{axis}[omrecord, width=0.78\linewidth, height=6.2cm,
  xmin=45, xmax=79, ymin=0, ymax=5,
  xlabel={\ce{SiO2} (\unit{\percent})}, ylabel={\ce{K2O} (\unit{\percent})},
  legend style={font=\scriptsize, at={(0.02,0.98)}, anchor=north west},
  tick label style={font=\scriptsize}, label style={font=\small}]
\addplot[only marks, mark=*, mark size=0.55pt, blue!70!black, opacity=0.6] table[x=SiO2, y=K2O] {figdata/out/bachelor-2-series-cath-th.dat};
\addlegendentry{IJsland}
\addplot[only marks, mark=*, mark size=0.55pt, red!75!black, opacity=0.6] table[x=SiO2, y=K2O] {figdata/out/bachelor-2-series-cath-ca.dat};
\addlegendentry{Cascaden}
\addplot[thick, black, forget plot] table[x=x, y=y_low] {figdata/out/bachelor-2-series-k2o-lines.dat};
\addplot[thick, black, forget plot] table[x=x, y=y_high] {figdata/out/bachelor-2-series-k2o-lines.dat};
\draw[thick] (axis cs:48,0) -- (axis cs:48,0.3);
\node[font=\scriptsize, fill=white, inner sep=1pt] at (axis cs:60,3.7) {hoog-K};
\node[font=\scriptsize, fill=white, inner sep=1pt] at (axis cs:76.5,2.2) {middel-K};
\node[font=\scriptsize, fill=white, inner sep=1pt] at (axis cs:60,0.35) {laag-K};
\end{axis}
\end{tikzpicture}
\par\smallskip{\small Kali tegen silica voor de IJslandse en Cascadische gesteenten, met
de lijnen door (48, 0.3)--(68, 1.2) en (48, 1.2)--(68, 3.1) die laag-K-,
middel-K- en hoog-K-typen scheiden (naar Peccerillo en Taylor, 1976, zoals overgenomen in het
internationale schema), hier voorbij 68 doorgetrokken.}
\end{omfigure}
% ledger: bgs:k2o-sio2, data:b4-cath

Waarom verschilt de oxidatie? Kalkalkalische \omterm{def:b2:magma-emplacement:magma}{magma's} ontstaan boven subductiezones, waar
de ontwaterende plaat geoxideerde species aan de mantelwig levert; de \omterm{def:b2:magma-emplacement:magma}{magma's}
kristalliseren magnetiet, en hun geëvolueerde leden hoornblende of biotiet. Tholeïtische
\omterm{def:b2:magma-emplacement:magma}{magma's} ontstaan bij ruggen, waar de mantel door decompressie smelt
(\cref{ch:b2:mantle-melting}), en zijn sterker gereduceerd.
% ledger: vp21:series

\section{Reeksen en omgevingen}\label{sec:b2:magmatic-series:settings}

Elke omgeving van de platentektoniek maakt haar eigen \omterm{def:b2:magma-emplacement:magma}{magma's}, omdat elke de
mantel op haar eigen manier doet smelten: door decompressie onder ruggen en hotspots, door water toe te voegen
boven subductiezones (\cref{ch:b2:mantle-melting}).

\emph{Mid-oceanische ruggen} brengen tholeïtische basalten voort, de MORB van
\cref{ch:b2:elements-in-the-earth}, arm aan \omterm{def:b2:elements-in-the-earth:partition}{incompatibele elementen} en met de
verarmde isotopensignatuur van hun bron. IJsland, een rug boven een hotspot, voegt
daar zijn eigen tholeïeten aan toe, met een volledige reeks tot rhyoliet.
% ledger: vp21:series

\emph{Oceaaneilanden} groeien in fasen. Een Hawaïaanse vulkaan brengt alkalische lava's voort zolang hij
jong is, en bouwt dan zijn schild van tholeïtisch basalt, naar schatting 80 tot
\qty{95}{\percent} van zijn uiteindelijke volume, binnen zijn eerste 1 tot 2 miljoen jaar; terwijl hij
van de hotspot wegdrijft, worden zijn lava's weer alkalisch, met als slot een lange staart van
zeldzame alkalische erupties. Minder smelten geeft alkalischere \omterm{def:b2:magma-emplacement:magma}{magma's}: de alkaliën worden, omdat ze
incompatibel zijn, in een kleine \omterm{def:b2:mantle-melting:fraction}{smeltfractie} minder verdund. Isotopisch bemonsteren
\omterm{def:b2:elements-in-the-earth:morb-oib}{oceaaneilandbasalten} de verrijkte reservoirs van \cref{ch:b2:isotope-geochemistry}.
% ledger: usgs14:hawaii-stages

\emph{Vulkanische bogen} brengen kalkalkalische reeksen voort. Hun \omterm{def:b2:elements-in-the-earth:trace}{sporenelementen} dragen een
kenmerkend stempel (de figuur hieronder): genormaliseerd op de \omterm{def:b2:elements-in-the-earth:bse}{primitieve mantel} zijn de
boogbasalten verrijkt in Rb, Ba, Th, U, Pb en Sr, elementen die de plaat aan de
wig toevoegt, en verarmd in niobium en tantaal ten opzichte van hun buren. De mediane
boogbasalt van de dataset heeft $\mathrm{Nb}/\mathrm{La} = 0.49$, tegen 0.97 voor MORB
en 1.23 voor \omterm{def:b2:elements-in-the-earth:morb-oib}{oceaaneilandbasalten}.
% ledger: data:b4-basalt-settings

\begin{omfigure}
\centering
\begin{tikzpicture}
\begin{axis}[omrecord, width=0.86\linewidth, height=6.4cm, ymode=log,
  xmin=0.5, xmax=22.5, ymin=1, ymax=300,
  xtick={1,...,22}, xticklabels={Rb,Ba,Th,U,Nb,Ta,La,Ce,Pb,Sr,Nd,Zr,Hf,Sm,Eu,Ti,Gd,Dy,Y,Er,Yb,Lu},
  ylabel={monster / primitieve mantel},
  legend style={font=\scriptsize, at={(0.98,0.98)}, anchor=north east},
  tick label style={font=\scriptsize}, label style={font=\small}]
\addplot[thick, mark=square*, mark size=1.4pt, orange!85!black] table[x=idx, y=oib] {figdata/out/bachelor-2-series-spider.dat};
\addlegendentry{oceaaneilandbasalt}
\addplot[thick, mark=*, mark size=1.4pt, blue!70!black] table[x=idx, y=morb] {figdata/out/bachelor-2-series-spider.dat};
\addlegendentry{MORB}
\addplot[thick, mark=triangle*, mark size=1.7pt, red!75!black] table[x=idx, y=iab] {figdata/out/bachelor-2-series-spider.dat};
\addlegendentry{eilandboogbasalt}
\end{axis}
\end{tikzpicture}
\par\smallskip{\small Op de \omterm{def:b2:elements-in-the-earth:bse}{primitieve mantel} genormaliseerde mediane samenstellingen van 241 MORB, 259
\omterm{def:b2:elements-in-the-earth:morb-oib}{oceaaneilandbasalten} en 256 eilandboogbasalten (de compilatie van Vermeesch, 2006;
\omterm{def:b2:elements-in-the-earth:bse}{primitieve mantel} van McDonough en Sun, 1995). De boogbasalten tonen het dal bij Nb
en Ta en pieken bij Ba, Pb en Sr; de \omterm{def:b2:elements-in-the-earth:morb-oib}{oceaaneilandbasalten} zijn verrijkt in alle
sterkst \omterm{def:b2:elements-in-the-earth:partition}{incompatibele elementen}.}
\end{omfigure}
% ledger: data:b4-basalt-settings, ms95:pm-trace

Binnen continenten ontmoet basaltisch \omterm{def:b2:magma-emplacement:magma}{magma} een dikke korst die het kan verhitten en doen smelten,
en het resultaat kan een vulkanisme met twee modi zijn.

\begin{definition}[Bimodaal vulkanisme]\label{def:b2:magmatic-series:bimodal}
\emph{Bimodaal vulkanisme}\index{bimodaal vulkanisme} is de uitbarsting, in één provincie,
van twee groepen \omterm{def:b2:magma-emplacement:magma}{magma's}, basaltisch en rhyolitisch, met weinig tussenliggende samenstellingen
ertussen.
\end{definition}

Het vulkanisme langs het spoor van de Yellowstone-hotspot is bimodaal: op de oostelijke
Snake River Plain worden oudere rhyolieten met zo'n \qty{75}{\percent} silica bedekt door
jongere basalten met zo'n \qty{50}{\percent}. De leemte ertussen is de aanwijzing: de
differentiatie van een basalt alleen zou tussenliggende lava's hebben achtergelaten, en een leemte
wijst op een tweede bron voor de rhyolieten, zoals de korst die de basalt verhit.
% ledger: yvo23:bimodal

\emph{Grote stollingsprovincies}, besproken in Boek~3, brengen basalt voort in volumes
die geen andere omgeving evenaart; hun reeks wordt met dezelfde diagrammen gelezen.

\section{Discriminatiediagrammen}\label{sec:b2:magmatic-series:discrimination}

De patronen van de vorige paragraaf nodigen uit om een omgeving rechtstreeks uit een analyse af te lezen,
en petrologen hebben daarvoor veel diagrammen getekend.

\begin{definition}[Discriminatiediagram]\label{def:b2:magmatic-series:discrimination}
Een \emph{discriminatiediagram}\index{discriminatiediagram} is een diagram van
chemische variabelen, meestal immobiele \omterm{def:b2:elements-in-the-earth:trace}{sporenelementen}, verdeeld in velden die zijn getekend
rond de samenstellingen van gesteenten uit bekende omgevingen, en gebruikt om een omgeving toe te kennen aan een
gesteente van onbekende herkomst.
\end{definition}

De klassieker is de Ti--Zr--Y-driehoek van Pearce en Cann (1973), getekend voor basalten. Hij
gebruikt elementen die door verwering weinig worden verplaatst, zodat hij kan worden toegepast op oude,
gemetamorfoseerde lava's waarvan de omgeving verloren is. De figuur hieronder test hem op de 538 basalten van de
compilatie die alle drie de elementen hebben. De meeste \omterm{def:b2:elements-in-the-earth:morb-oib}{oceaaneilandbasalten} vallen in het
intraplaatveld D, 152 van de 173, en de meeste MORB in veld B, 142 van de 164; maar veld B wordt
gedeeld met boogbasalten, en 101 van de 201 boogbasalten vallen er ook in.
% ledger: geoplotr:tizry, data:b4-basalt-settings

\begin{omfigure}
\centering
\begin{tikzpicture}
\begin{axis}[width=0.62\linewidth, height=0.55\linewidth, axis equal image,
  hide axis, xmin=-0.04, xmax=1.04, ymin=-0.06, ymax=0.92, clip=false,
  legend style={font=\scriptsize, at={(1.0,0.95)}, anchor=north west}]
\addplot[only marks, mark=*, mark size=0.7pt, blue!70!black, opacity=0.7] table[x=x, y=y] {figdata/out/bachelor-2-series-tizry-morb.dat};
\addlegendentry{MORB}
\addplot[only marks, mark=square*, mark size=0.7pt, orange!85!black, opacity=0.7] table[x=x, y=y] {figdata/out/bachelor-2-series-tizry-oib.dat};
\addlegendentry{oceaaneilandbasalt}
\addplot[only marks, mark=triangle*, mark size=0.9pt, red!75!black, opacity=0.7] table[x=x, y=y] {figdata/out/bachelor-2-series-tizry-iab.dat};
\addlegendentry{eilandboogbasalt}
\foreach \f in {A,B,C,D}
  \addplot[thick, black, forget plot] table[x=x, y=y] {figdata/out/bachelor-2-series-tizry-\f.dat};
\draw[thick] (axis cs:0,0) -- (axis cs:1,0) -- (axis cs:0.5,0.8660254) -- cycle;
\node[anchor=south] at (axis cs:0.5,0.866) {Ti/100};
\node[anchor=north east] at (axis cs:0,0) {Zr};
\node[anchor=north west] at (axis cs:1,0) {Y $\times$ 3};
\draw[thin] (axis cs:0.33,0.36) -- (axis cs:0.22,0.46) node[anchor=south east, font=\small\bfseries, inner sep=1pt] {D};
\draw[thin] (axis cs:0.36,0.12) -- (axis cs:0.24,0.06) node[anchor=north east, font=\small\bfseries, inner sep=1pt] {C};
\draw[thin] (axis cs:0.56,0.37) -- (axis cs:0.76,0.50) node[anchor=west, font=\small\bfseries, inner sep=1pt] {A};
\draw[thin] (axis cs:0.53,0.21) -- (axis cs:0.72,0.13) node[anchor=north west, font=\small\bfseries, inner sep=1pt] {B};
\end{axis}
\end{tikzpicture}
\par\smallskip{\small Het Ti--Zr--Y-diagram van Pearce en Cann (1973) met de basalten
van bekende omgeving. A eilandboogtholeïeten; B MORB, eilandboogtholeïeten en
kalkalkalische basalten; C kalkalkalische basalten; D intraplaatbasalten. Ti uit
\ce{TiO2}; alle drie in ppm.}
\end{omfigure}
% ledger: geoplotr:tizry, data:b4-basalt-settings

\begin{method}[Een discriminatiediagram gebruiken]\label{met:b2:magmatic-series:discrimination}
(1) Ga na of het gesteente van de soort is waarvoor het diagram is gemaakt, meestal een basalt, en
geen \omterm{def:b2:magmatic-differentiation:cumulate}{cumulaat} of geëvolueerde lava, waarvan de \omterm{def:b2:elements-in-the-earth:trace}{sporenelementen} evenzeer de kristallisatie weerspiegelen
als de bron. (2) Geef bij verweerde of gemetamorfoseerde gesteenten de voorkeur aan elementen die bij verwering immobiel zijn (Ti, Zr, Y, Nb, de zware
\omterm{def:b2:elements-in-the-earth:ree}{zeldzame aarden}). (3) Zet een reeks monsters uit, niet één:
een enkel punt dicht bij een grens beslist niets. (4) Lees het resultaat als een waarschijnlijkheid:
de velden zijn getekend rond de meeste, niet alle, van hun referentiemonsters, en sommige velden
worden gedeeld. (5) Bevestig met onafhankelijke aanwijzingen: andere diagrammen, isotopen, de
geologische context.
\end{method}

Stap (4) kan kwantitatief worden gemaakt (\cref{exo:b2:magmatic-series:8}). Als een basalt in veld B valt en de
drie omgevingen even vaak voorkwamen, is de kans dat het een MORB is maar ongeveer 0.62.
Statistische classificatoren die op zulke compilaties zijn getraind, erven dezelfde grens: een basalt
draagt de herinnering aan zijn bron en zijn smelten mee, en twee omgevingen kunnen vergelijkbare
\omterm{def:b2:magma-emplacement:magma}{magma's} maken.
% ledger: data:b4-basalt-settings

\begin{inthelab}[Een analyse van het totale gesteente controleren]
Voordat ze een gesteente in een van deze diagrammen uitzetten, controleren petrologen de analyse. De
hoofdoxiden, bijvoorbeeld gemeten met röntgenfluorescentie, moeten samen met de vluchtige bestanddelen
dicht bij 100 uitkomen: water en koolstofdioxide, rechtstreeks gemeten of samen als het
gloeiverlies, de massa die een poeder bij sterke verhitting verliest. Een groot
gloeiverlies wijst op een verweerd gesteente, waarvan de alkaliën verplaatst kunnen zijn; het schema van het
TAS-diagram verwerpt analyses met meer dan \qty{2}{\percent} water of
\qty{0.5}{\percent} koolstofdioxide. Analyses worden daarna herrekend tot 100 zonder
vluchtige bestanddelen, ijzer wordt uitgedrukt als totaal \ce{FeO} of verdeeld over ijzer(II) en ijzer(III),
en referentiematerialen die in dezelfde reeks zijn geanalyseerd, worden vergeleken met hun gecertificeerde
waarden (\cref{ch:b2:field-lab-methods-2}).
\end{inthelab}
% ledger: bgs:tas-use

\section{Oefeningen}

\begin{exercise}[$\star$]\label{exo:b2:magmatic-series:1}
Benoem de volgende vulkanische gesteenten in het TAS-diagram (\ce{SiO2}, \ce{Na2O}, \ce{K2O}
in \unit{\percent}, al vrij van vluchtige bestanddelen): (a) 49.60, 2.22, 0.51 (de Hawaïaanse
referentiebasalt BHVO-2); (b) 58.5, 3.9, 1.6; (c) 72.0, 4.0, 4.2; (d) 47.0, 4.5, 2.5.
Noem voor (d) de twee mogelijke gesteenten en zeg wat tussen beide beslist.
\end{exercise}
% ledger: usgs:bhvo2-major, bgs:tas

\begin{exercise}[$\star$]\label{exo:b2:magmatic-series:2}
Een lava geeft \ce{SiO2} 51.60, \ce{TiO2} 1.20, \ce{Al2O3} 16.10, \ce{FeO_T} 9.80,
\ce{MgO} 6.60, \ce{CaO} 8.20, \ce{Na2O} 3.50, \ce{K2O} 1.40, \ce{H2O} 1.40 en
\ce{CO2} \qty{0.20}{\percent}. (a) Ga na of de analyse aanvaardbaar is voor het TAS-diagram.
(b) Herreken haar zonder vluchtige bestanddelen. (c) Benoem het gesteente voor en na de
herberekening. (d) Is het laag-, middel- of hoog-K?
\end{exercise}
% ledger: bgs:tas-use, bgs:k2o-sio2

\begin{exercise}[$\star$]\label{exo:b2:magmatic-series:3}
Benoem in het \omterm{def:b2:magmatic-series:qapf}{QAPF-diagram}: (a) een grofkorrelig gesteente met \qty{25}{\percent} kwarts,
\qty{30}{\percent} alkaliveldspaat, \qty{35}{\percent} plagioklaas, \qty{8}{\percent}
biotiet en \qty{2}{\percent} hoornblende; (b) een met \qty{5}{\percent} kwarts,
\qty{5}{\percent} alkaliveldspaat, \qty{55}{\percent} plagioklaas en
\qty{35}{\percent} clinopyroxeen.
\end{exercise}
% ledger: bgs:qapf

\begin{exercise}[$\star$]\label{exo:b2:magmatic-series:4}
Deel in als laag-, middel- of hoog-K, met de lijnen van de \ce{K2O}--\ce{SiO2}-figuur:
(a) 50.0, 0.40; (b) 56.0, 1.50; (c) 62.0, 2.90; (d) 66.0, 1.00 (\ce{SiO2},
\ce{K2O}, \unit{\percent}).
\end{exercise}
% ledger: bgs:k2o-sio2

\begin{exercise}[$\star\star$]\label{exo:b2:magmatic-series:5}
(a) Bereken $A$, $F$ en $M$ voor een basalt met \ce{FeO_T} 14.5, \ce{MgO} 4.5 en
\ce{Na2O + K2O} \qty{3.2}{\percent}, en voor een andesiet met 6.1, 3.8 en 5.05.
(b) De scheidingslijn gaat door $(F, A, M) = (54.5, 23.0, 22.5)$, $(52.5, 30.0,
17.5)$ en $(49.6, 35.0, 15.4)$, waarvan het eerste het ijzerrijkste punt is; de
tholeïtische kant is de kant van F. Bepaal aan welke kant elk gesteente valt. (c) Welk
mineraal houdt, door vroeg te kristalliseren, de andesiet waar hij is?
\end{exercise}
% ledger: geoplotr:afm, vp21:series

\begin{exercise}[$\star\star$]\label{exo:b2:magmatic-series:6}
De mediane MORB van de compilatie heeft \qty{4}{ppm} Nb en \qty{4.12}{ppm} La, de
mediane \omterm{def:b2:elements-in-the-earth:morb-oib}{oceaaneilandbasalt} 38 en 30.8, de mediane boogbasalt 3 en 6.1. Bereken met de
waarden van de \omterm{def:b2:elements-in-the-earth:bse}{primitieve mantel} Nb 0.658 en La \qty{0.648}{ppm} de genormaliseerde waarden
en $\mathrm{Nb_N}/\mathrm{La_N}$ voor elk. Welke heeft een niobiumanomalie, en waarom?
\end{exercise}
% ledger: data:b4-basalt-settings, ms95:pm-trace

\begin{exercise}[$\star\star$]\label{exo:b2:magmatic-series:7}
Ti is $0.5993$ van de massa van \ce{TiO2}. Bereken de coördinaten in het Ti--Zr--Y-diagram
(Ti/100, Zr, $3\mathrm Y$, herrekend tot 100) van drie basalten: (a) \ce{TiO2}
\qty{1.6}{\percent}, Zr 110, Y \qty{32}{ppm}; (b) 2.8, 220, 29; (c) 0.8, 50, 20. Plaats
ze in de figuur en zeg welke conclusie elke plaatsing toelaat.
\end{exercise}
% ledger: geoplotr:tizry

\begin{exercise}[$\star\star$]\label{exo:b2:magmatic-series:8}
In de Ti--Zr--Y-test van dit hoofdstuk bevat veld B 142 van de 164 MORB, 4 van de 173
\omterm{def:b2:elements-in-the-earth:morb-oib}{oceaaneilandbasalten} en 101 van de 201 boogbasalten. (a) Bepaal welke fractie van elke groep
in B valt. (b) Bereken, als de drie omgevingen even vaak voorkwamen, de kans dat een
basalt in B een MORB is. (c) Wat valt er op grond van veld B alleen te zeggen?
\end{exercise}
% ledger: data:b4-basalt-settings

\begin{exercise}[$\star\star$]\label{exo:b2:magmatic-series:9}
Een rhyoliet heeft \ce{Na2O} 5.5, \ce{K2O} 4.5 en \ce{Al2O3} \qty{10.5}{\percent} ($M$ =
61.98, 94.20 en \qty{101.96}{g/mol}). Bereken zijn moleculaire $(\ce{Na2O +
K2O})/\ce{Al2O3}$. (a) Is hij peralkalisch? (b) Alkaliveldspaat bevat één
alkalioxide per \ce{Al2O3}: bepaal welke fractie van de alkaliën er niet in kan.
\end{exercise}
% ledger: bgs:tas-use

\begin{exercise}[$\star\star\star$]\label{exo:b2:magmatic-series:10}
Neem in het model van de propositie over log-ratiotrends $\lambda_A = 0.4$ en $\lambda_M =
0.8$. (a) Bereken de helling van de trend voor $\lambda_F = 0.5$ en voor $\lambda_F =
0.7$. (b) De IJslandse gesteenten hebben een helling van $-1.10$ en die van de Cascaden
$-0.72$: bepaal welke $\lambda_F$ elk impliceert. (c) Welke reeks verliest haar ijzer sneller,
en welk mineraal verklaart dat? (d) Wat gebeurt er met $F/M$ langs elke trend?
\end{exercise}
% ledger: vp21:model

\begin{exercise}[$\star\star\star$]\label{exo:b2:magmatic-series:11}
Een vereenvoudigde norm. Werk in mol oxide. Zet al het \ce{K2O} in orthoklaas
(\ce{K2O.Al2O3.6SiO2}), al het \ce{Na2O} in albiet (\ce{Na2O.Al2O3.6SiO2}), het
resterende \ce{Al2O3} met evenveel \ce{CaO} in anorthiet (\ce{CaO.Al2O3.2SiO2}), het
resterende \ce{CaO} met evenveel $(\ce{Mg,Fe})\ce{O}$ in diopsied
(\ce{CaO.(Mg,Fe)O.2SiO2}) en de rest van het $(\ce{Mg,Fe})\ce{O}$ in hypersteen
(\ce{(Mg,Fe)O.SiO2}). Is er te weinig silica, zet dan hypersteen om in olivijn
(\ce{2(Mg,Fe)O.SiO2}); is er nog steeds te weinig, zet dan albiet om in nefelien
(\ce{Na2O.Al2O3.2SiO2}). Pas dit toe op (a) BHVO-2: \ce{SiO2} 49.60, \ce{Al2O3} 13.44,
\ce{FeO_T} 11.15, \ce{MgO} 7.257, \ce{CaO} 11.40, \ce{Na2O} 2.219, \ce{K2O}
\qty{0.513}{\percent}; (b) een basalt met 45.5, 15.0, 11.5, 8.5, 10.0, 3.8 en 1.5 (molaire
massa's \ce{SiO2} 60.08, \ce{Al2O3} 101.96, \ce{FeO} 71.84, \ce{MgO} 40.30, \ce{CaO}
56.08, \ce{Na2O} 61.98, \ce{K2O} \qty{94.20}{g/mol}). Ga na of elk verzadigd is aan silica. Bepaal
tot welke reeks elke basalt behoort.
\end{exercise}
% ledger: usgs:bhvo2-major, bgs:series-terms

\begin{exercise}[$\star\star\star$]\label{exo:b2:magmatic-series:12}
In de compilatie geldt $\mathrm{Nb}/\mathrm{La} < 0.7$ voor 116 van de 141 boogbasalten,
21 van de 101 MORB en 3 van de 137 \omterm{def:b2:elements-in-the-earth:morb-oib}{oceaaneilandbasalten} waarin beide zijn geanalyseerd. (a) Bereken de
drie fracties. (b) Bereken, bij gelijke aandelen van de drie omgevingen, de
kans dat een basalt met $\mathrm{Nb}/\mathrm{La} < 0.7$ een boogbasalt is. (c)
Herhaal dit voor een gebied waar \qty{60}{\percent} van de basalten MORB is, \qty{30}{\percent}
\omterm{def:b2:elements-in-the-earth:morb-oib}{oceaaneilandbasalt} en \qty{10}{\percent} boogbasalt. (d) Welke les houdt (c)
in voor het gebruik van discriminatiediagrammen?
\end{exercise}
% ledger: data:b4-basalt-settings

\section{Probleem: Eén vulkaan, drie lava's}

\begin{problem}[{Weekendprobleem --- van drie analyses naar de omgeving van een vulkaan en
de kristallen die hij achterliet}]
\label{pb:b2:magmatic-series:1}
Drie lava's van een (verzonnen) stratovulkaan, van zijn oudste flankstromen tot zijn jongste
koepel, geven (\unit{\percent}):

\smallskip\noindent{\small
\begin{tabular}{@{}lrrrrrrrrrr@{}}
 & \ce{SiO2} & \ce{TiO2} & \ce{Al2O3} & \ce{FeO_T} & \ce{MgO} & \ce{CaO} & \ce{Na2O} & \ce{K2O} & \ce{H2O} & \ce{CO2}\\
lava 1 & 50.4 & 0.95 & 17.6 & 8.6 & 7.9 & 10.2 & 2.9 & 0.62 & 0.70 & 0.10\\
lava 2 & 58.6 & 0.80 & 17.4 & 6.1 & 3.8 & 6.9 & 3.9 & 1.15 & 1.10 & 0.05\\
lava 3 & 65.3 & 0.55 & 16.2 & 4.0 & 1.6 & 4.3 & 4.4 & 1.85 & 1.20 & 0.05\\
\end{tabular}}

\smallskip\noindent en in ppm: Rb 10, 20, 32; Th 1.1, 2.1, 3.4; Nb 2.0, 3.6, 5.6;
La 7, 13, 20; Sr 520, 480, 380; Zr 70, 125, 190; Y 18, 22, 24. \omterm{def:b2:elements-in-the-earth:bse}{Primitieve mantel}: Nb
0.658, La 0.648, Th \qty{0.0795}{ppm}.

\noindent\textbf{Deel I --- De lava's benoemen.}
\begin{enumerate}
\item Tel elke analyse op. Zijn de drie aanvaardbaar voor het TAS-diagram?
\item Herreken lava 1 zonder vluchtige bestanddelen: geef de factor, \ce{SiO2} en
\ce{Na2O + K2O}.
\item Doe hetzelfde voor lava 2 en 3.
\item Benoem de drie lava's in het TAS-diagram.
\item Zouden de namen zonder de herberekening veranderen?
\item Een grofkorrelig gesteente met de samenstelling van lava 2 kristalliseerde in de wortels van de
vulkaan. Zeg welk schema het benoemt, en in welk veld het te verwachten is.
\end{enumerate}

\noindent\textbf{Deel II --- Welke reeks?}
\begin{enumerate}[resume]
\item Deel elke lava in als laag-, middel- of hoog-K.
\item Bereken $A$, $F$, $M$ voor elk (zonder vluchtige bestanddelen).
\item De scheidingslijn gaat door $(F, A, M) = (42.8, 17.1, 40.1)$, $(46.9, 18.1,
35.0)$, $(52.5, 30.0, 17.5)$, $(49.6, 35.0, 15.4)$, $(40.8, 50.0, 9.2)$ en $(37.9,
55.0, 7.1)$, met het tholeïtische veld aan de kant van F. Interpoleer $F$ lineair in
$A$ tussen naburige punten en bepaal aan welke kant elke lava valt.
\item Bereken \ce{FeO_T}/\ce{MgO} voor elke lava. Is er een ijzerverrijking?
\item Zeg tot welke reeks de vulkaan behoort. Welk mineraal en welke omstandigheden
verklaren dat?
\end{enumerate}

\noindent\textbf{Deel III --- Welke omgeving?}
\begin{enumerate}[resume]
\item Normaliseer Nb, La en Th van lava 1 op de \omterm{def:b2:elements-in-the-earth:bse}{primitieve mantel}.
\item Bereken $\mathrm{Nb_N}/\mathrm{La_N}$ voor lava 1 en vergelijk met de medianen
van \cref{exo:b2:magmatic-series:6}.
\item Bereken $\mathrm{Th_N}/\mathrm{Nb_N}$. Wat voerde het thorium aan?
\item Bereken de Ti--Zr--Y-coördinaten van lava 1 (\ce{TiO2} van de analyse zonder vluchtige
bestanddelen). Wat zegt het veld ervan? Waarom lava 2 en 3 niet uitzetten?
\item Concludeer: waar is deze vulkaan gegroeid?
\end{enumerate}

\noindent\textbf{Deel IV --- Hoeveel kristalliseerde er?}
\begin{enumerate}[resume]
\item Kali gaat in geen van de vroege mineralen. Bereken uit het \ce{K2O} zonder vluchtige bestanddelen van
lava 1 en 2 de fractie $F$ vloeistof die over was toen lava 2 ontstond, als die
door \omterm{def:b2:magmatic-differentiation:crystallisation}{fractionele kristallisatie} uit lava 1 voortkomt.
\item Doe hetzelfde met rubidium en met thorium.
\item Zirkonium heeft een globale $D \approx 0.1$. Gebruik de wet van Rayleigh om er $F$
uit te schatten.
\item Strontium daalt van lava 2 naar lava 3. Bereken $F$ uit kali tussen lava 2 en
3, en schat daarmee de globale $D$ van Sr over dat interval. Welk mineraal is
verantwoordelijk?
\item Bepaal welke fractie van de vloeistof van lava 1 tot lava 3 kristalliseerde.
\item Zou de schatting met kali te hoog of te laag zijn als het \omterm{def:b2:magma-emplacement:magma}{magma} kalirijke granitische korst had
geassimileerd?
\item Neem het gemiddelde van de vier schattingen van Deel IV voor lava 1 tot 2 en geef de massa
kristallen die per \qty{100}{kg} basalt achterbleef om de andesiet te maken, het kerngetal
van dit probleem.
\end{enumerate}
\end{problem}
% ledger: bgs:tas-use, bgs:tas, bgs:k2o-sio2, geoplotr:afm, geoplotr:tizry, ms95:pm-trace
